% ============================================================
%  Chapitre 4 : Réalisation
% ============================================================

% ---- Page de titre du chapitre ----
\clearpage
\thispagestyle{chaptertitlepage}
\vspace*{\fill}
\begin{center}
\rule{\textwidth}{0.4pt}\\[0.8cm]
{\fontsize{24}{29}\selectfont\bfseries Chapitre 4}\\[0.6cm]
{\fontsize{18}{22}\selectfont Réalisation}\\[0.8cm]
\rule{\textwidth}{0.4pt}
\end{center}
\vspace*{\fill}
\clearpage

% ---- Contenu du chapitre ----
\chapter{Réalisation}

\section{Introduction}

La phase de réalisation a consisté à implémenter l'ensemble des fonctionnalités spécifiées dans l'étude fonctionnelle et technique, et modélisées dans les diagrammes UML. Le développement a suivi une approche itérative, en commençant par la couche backend (API et base de données), puis la couche frontend (interface utilisateur).

L'environnement de développement comprend~:
\begin{itemize}
    \item \textbf{IDE}~: Visual Studio Code avec les extensions C\#, TypeScript, Tailwind CSS IntelliSense
    \item \textbf{Backend}~: \texttt{dotnet run} (port HTTPS 7091, HTTP 5057)
    \item \textbf{Frontend}~: \texttt{npm run dev} (port 8080, proxy vers le backend)
    \item \textbf{Base de données}~: SQL Server LocalDB (\texttt{(localdb)\textbackslash MSSQLLocalDB}, base \texttt{CabinetD})
\end{itemize}

% ============================================================
\section{Sécurité~: Authentification et Contrôle d'accès}

La sécurité de l'application repose sur le standard \textbf{JSON Web Token (JWT)}, un mécanisme d'authentification \emph{stateless} largement adopté dans les architectures REST modernes.

\subsection{Configuration JWT côté backend}

Dans le fichier \texttt{Program.cs}, le middleware d'authentification ASP.NET Core est configuré au démarrage~:

\begin{lstlisting}[language={[Sharp]C}, caption={Configuration JWT dans Program.cs}]
builder.Services.AddAuthentication(JwtBearerDefaults.AuthenticationScheme)
    .AddJwtBearer(options => {
        options.TokenValidationParameters = new TokenValidationParameters {
            ValidateIssuer = true,
            ValidateAudience = true,
            ValidateLifetime = true,
            ValidateIssuerSigningKey = true,
            ValidIssuer = configuration["Jwt:Issuer"],
            ValidAudience = configuration["Jwt:Audience"],
            IssuerSigningKey = new SymmetricSecurityKey(
                Encoding.UTF8.GetBytes(configuration["Jwt:Key"]))
        };
    });
\end{lstlisting}

Les paramètres JWT sont stockés dans \texttt{appsettings.json}~:
\begin{itemize}
    \item \texttt{Jwt:Key}~: clé secrète HMAC-SHA256 pour la signature du token
    \item \texttt{Jwt:Issuer}~: identifiant de l'émetteur (le backend)
    \item \texttt{Jwt:Audience}~: identifiant de l'audience (le frontend)
\end{itemize}

Le \texttt{AuthController} expose un unique endpoint \texttt{POST /api/Auth/login} qui~:
\begin{enumerate}
    \item Reçoit un \texttt{LoginDto} contenant \texttt{username} et \texttt{password}.
    \item Recherche l'utilisateur dans la table \texttt{Usercmim} par \texttt{username}.
    \item Compare le mot de passe fourni avec celui stocké en base.
    \item Si valide, crée un \texttt{SecurityTokenDescriptor} avec un claim \texttt{ClaimTypes.Name}, une expiration de 1 heure, et une signature HMAC-SHA256.
    \item Retourne le token JWT sérialisé et les informations utilisateur.
\end{enumerate}

\subsection{Protection côté frontend}

Le frontend utilise un \textbf{pattern Provider/Context} (React Context API) pour gérer l'authentification de manière globale~:

\begin{itemize}
    \item \textbf{AuthProvider} (\texttt{use-auth.tsx})~: enveloppe toute l'application et expose via le contexte~:
    \begin{itemize}
        \item \texttt{user}~: l'utilisateur connecté (ou null)
        \item \texttt{isAuthenticated}~: booléen dérivé de la présence du token
        \item \texttt{login(email, password)}~: appelle l'API, stocke le token dans \texttt{localStorage}
        \item \texttt{logout()}~: supprime le token, redirige vers la page de connexion
        \item \texttt{updateUser(userData)}~: met à jour les données utilisateur locales
    \end{itemize}
    
    \item \textbf{ProtectedRoute}~: composant wrapper qui vérifie \texttt{isAuthenticated}. Si l'utilisateur n'est pas connecté, redirection automatique vers \texttt{/}.
    
    \item \textbf{ApiService.getAuthHeaders()}~: méthode qui injecte automatiquement l'en-tête \texttt{Authorization: Bearer <token>} dans chaque requête API.
\end{itemize}

\subsection{Swagger avec authentification}

En environnement de développement, Swagger UI est activé avec un bouton \og Authorize \fg{} permettant de saisir le token JWT. La configuration inclut un \texttt{SecurityDefinition} de type Bearer et un \texttt{SecurityRequirement} global, permettant de tester tous les endpoints protégés directement depuis le navigateur.

% ============================================================
\section{Gestion des requêtes et communication client-serveur}

\subsection{Service API centralisé}

Toute la communication entre le frontend et le backend est centralisée dans la classe \textbf{ApiService} (\texttt{src/lib/api.ts}). Cette classe fournit une abstraction de haut niveau sur l'API Fetch~:

\begin{itemize}
    \item \textbf{Authentification}~: \texttt{login()}, \texttt{logout()}, \texttt{isAuthenticated()}, \texttt{getToken()}, \texttt{getAuthHeaders()}
    \item \textbf{Patients}~: \texttt{getPatients()}, \texttt{createPatient()}, \texttt{updatePatient()}, \texttt{deletePatient()}, \texttt{getPatientStats()}, \texttt{getPatientDossier()}
    \item \textbf{Rendez-vous}~: \texttt{getAppointments()}, \texttt{createAppointment()}, \texttt{updateAppointment()}, \texttt{deleteAppointment()}, \texttt{getAvailableSlots()}
    \item \textbf{Profil}~: \texttt{getProfile()}, \texttt{updateProfile()}, \texttt{changePassword()}
\end{itemize}

Chaque méthode gère automatiquement l'injection des headers d'authentification, le \texttt{Content-Type: application/json}, le parsing JSON et la gestion d'erreurs.

\subsection{Gestion du cache avec React Query}

Le frontend utilise \textbf{TanStack React Query} pour optimiser les interactions avec l'API~:

\begin{itemize}
    \item \textbf{Queries} (\texttt{useQuery})~: les données sont mises en cache avec des clés uniques (ex~: \texttt{["patients"]}, \texttt{["appointments"]}). Le cache évite les requêtes réseau redondantes.
    \item \textbf{Mutations} (\texttt{useMutation})~: les opérations d'écriture invalident automatiquement les caches associés via \texttt{queryClient.invalidateQueries()} après succès, déclenchent des notifications toast et gèrent les états de chargement.
\end{itemize}

\subsection{Détection des conflits de rendez-vous}

La détection des conflits est implémentée côté serveur dans le \texttt{RendezVousController}~:

\begin{itemize}
    \item \textbf{Conflit médecin}~: vérifie qu'aucun rendez-vous actif (statut $\neq$ \og canceled \fg) n'existe pour le même médecin, à la même date et heure. Si conflit $\rightarrow$ \texttt{409 Conflict}.
    \item \textbf{Conflit cabinet}~: vérifie qu'aucun rendez-vous actif n'existe pour le même cabinet, à la même date et heure. Si conflit $\rightarrow$ \texttt{409 Conflict}.
\end{itemize}

Cette double vérification garantit qu'un médecin ne peut pas avoir deux consultations simultanées et qu'un cabinet ne peut pas être utilisé par deux médecins en même temps.

% ============================================================
\section{Notifications et retour utilisateur}

L'application implémente un système de notifications visuelles basé sur les composants \textbf{Toast} de shadcn/ui~:

\begin{itemize}
    \item \textbf{Notifications toast}~: chaque action déclenche une notification éphémère (vert = succès, rouge = erreur, bleu = information).
    \item \textbf{Panneau de notifications} (\texttt{NotificationsPanel.tsx})~: accessible via l'icône cloche dans la barre de navigation, il affiche l'historique avec horodatage relatif et un badge de compteur.
\end{itemize}

% ============================================================
\section{Interface d'administration --- Pages de l'application}

\subsection{Page de connexion (Login)}

 \begin{figure}[H]
 \centering
 \includegraphics[width=0.8\textwidth]{screenshot_login.png}
 \caption{Capture d'écran~: Page de connexion}
 \label{fig:screenshot_login}
 \end{figure}

La page de connexion (\texttt{Login.tsx}) présente~:
\begin{itemize}
    \item Un arrière-plan avec dégradé bleu-violet et effet glassmorphism
    \item Une carte centrale avec logo \og DentalCare \fg, champ email, champ mot de passe
    \item Bouton \og Sign in \fg{} avec animation de chargement (spinner)
    \item Gestion d'erreurs via notification toast rouge
    \item Redirection automatique vers \texttt{/dashboard} après succès
\end{itemize}

\subsection{Tableau de bord (Dashboard)}

 \begin{figure}[H]
 \centering
 \includegraphics[width=0.85\textwidth]{screenshot_dashboard.png}
 \caption{Capture d'écran~: Tableau de bord}
 \label{fig:screenshot_dashboard}
 \end{figure}

Le Dashboard (\texttt{Dashboard.tsx}) affiche~:

\textbf{4 cartes de statistiques~:}
\begin{table}[H]
\centering
\begin{tabular}{@{} l l l l @{}}
\toprule
\textbf{Carte} & \textbf{Icône} & \textbf{Donnée} & \textbf{Couleur} \\
\midrule
Total Appointments & Calendar & Total RDV en base & Bleu \\
Today's Appointments & Clock & RDV du jour & Vert \\
Total Patients & Users & Patients actifs & Violet \\
Available Slots & CheckCircle & Créneaux libres & Orange \\
\bottomrule
\end{tabular}
\end{table}

\textbf{Section \og Recent Appointments \fg~:} les 6 premiers RDV du jour avec nom, type de soin, heure et badge de statut.

\textbf{Section \og Quick Actions \fg~:} quatre boutons d'accès rapide vers les pages principales.

\subsection{Gestion des patients (Patients)}

 \begin{figure}[H]
 \centering
 \includegraphics[width=0.85\textwidth]{screenshot_patients.png}
 \caption{Capture d'écran~: Gestion des patients}
 \label{fig:screenshot_patients}
 \end{figure}

La page Patients (\texttt{Patients.tsx}) offre~:
\begin{itemize}
    \item Bouton \og + New Patient \fg{} ouvrant une modale de création
    \item Barre de recherche avec filtrage en temps réel par nom, téléphone, email
    \item Filtres par matricule et par statut (Active, New, Inactive)
    \item Tableau avec colonnes~: Name, Age, Gender, Last Visit, Status, Phone, Email, Actions
    \item Actions par ligne~: Dossier, Éditer, Supprimer
    \item Modale de création/édition avec 16 champs
\end{itemize}

\subsection{Gestion des rendez-vous (Appointments)}

 \begin{figure}[H]
 \centering
 \includegraphics[width=0.85\textwidth]{screenshot_appointments.png}
 \caption{Capture d'écran~: Gestion des rendez-vous}
 \label{fig:screenshot_appointments}
 \end{figure}

La page Appointments (\texttt{Appointments.tsx}) offre~:
\begin{itemize}
    \item Bouton \og + New Appointment \fg{} avec modale contenant~:
    \begin{itemize}
        \item Sélecteur de patient (Combobox recherchable avec nom + matricule)
        \item Champs~: médecin, cabinet, date, heure, durée (15/30/45/60 min), statut, nature du soin, notes
    \end{itemize}
    \item Filtre par statut (Confirmed, Pending, Canceled)
    \item Tableau~: Patient, Doctor, Cabinet, Date \& Time, Type, Duration, Status, Actions
    \item Badges colorés selon le statut (vert/jaune/rouge)
\end{itemize}

\subsection{Calendrier hebdomadaire (Calendar)}

 \begin{figure}[H]
 \centering
 \includegraphics[width=0.85\textwidth]{screenshot_calendar.png}
 \caption{Capture d'écran~: Calendrier hebdomadaire}
 \label{fig:screenshot_calendar}
 \end{figure}

La page Calendar (\texttt{Calendar.tsx}) présente~:
\begin{itemize}
    \item Navigation semaine par semaine avec boutons précédent/suivant
    \item Grille 8 colonnes (heure + 7 jours) $\times$ 20 lignes (08:00 à 17:30, pas de 30 min)
    \item En-têtes des jours avec mise en surbrillance du jour courant
    \item Rendez-vous dans les cellules avec bordure latérale colorée selon le statut
    \item Indicateur rouge de l'heure courante sur la colonne du jour
    \item Panneau latéral \og Today's Schedule \fg{} avec les RDV du jour
    \item Panneau \og Quick Stats \fg{} (total du jour, premier créneau)
\end{itemize}

\subsection{Créneaux disponibles (Available Slots)}

 \begin{figure}[H]
 \centering
 \includegraphics[width=0.85\textwidth]{screenshot_slots.png}
 \caption{Capture d'écran~: Créneaux disponibles}
 \label{fig:screenshot_slots}
 \end{figure}

La page Available Slots (\texttt{AvailableSlots.tsx}) affiche~:
\begin{itemize}
    \item Navigation par jour avec boutons et sélecteur de date
    \item Panneau de configuration~: durée du créneau (15/30/45/60 min), statistiques (total, disponibles, réservés, taux d'utilisation)
    \item Grille responsive de créneaux~: vert = disponible (cliquable), gris = réservé
    \item Modale de réservation rapide au clic sur un créneau disponible
\end{itemize}

\subsection{Dossier patient (Patient Dossier)}

 \begin{figure}[H]
 \centering
 \includegraphics[width=0.85\textwidth]{screenshot_dossier.png}
 \caption{Capture d'écran~: Dossier patient}
 \label{fig:screenshot_dossier}
 \end{figure}

La page PatientDossier (\texttt{PatientDossier.tsx}) affiche~:
\begin{itemize}
    \item En-tête avec nom du patient, badge de genre, bouton retour et bouton \og Imprimer PDF \fg
    \item Carte d'informations~: date de naissance, dernière visite, prochain RDV, téléphone, email
    \item Tableau d'historique des rendez-vous (date, heure, docteur, nature du soin, statut)
    \item Export PDF via \texttt{window.print()} avec feuille de style \texttt{@media print} dédiée
\end{itemize}

\subsection{Profil utilisateur (Profile)}

 \begin{figure}[H]
 \centering
 \includegraphics[width=0.7\textwidth]{screenshot_profile.png}
 \caption{Capture d'écran~: Page de profil}
 \label{fig:screenshot_profile}
 \end{figure}

La page Profile (\texttt{Profile.tsx}) permet~:
\begin{itemize}
    \item Section \og Profile Information \fg~: email (lecture seule), nom (modifiable)
    \item Section \og Change Password \fg~: ancien mot de passe + nouveau mot de passe
    \item Données chargées via \texttt{GET /api/Profile/me}, modifications via \texttt{PUT}
\end{itemize}

\subsection{Navigation et mise en page}

\textbf{Barre latérale} (\texttt{AppSidebar.tsx})~:
\begin{itemize}
    \item Logo \og DentalCare Pro \fg{} en en-tête
    \item Menu~: Dashboard, Calendar, Appointments, Patients, Available Slots, Settings
    \item Mode rétractable (icône seule)
    \item Mise en surbrillance de la page active via \texttt{useLocation}
\end{itemize}

\textbf{Barre de navigation supérieure} (\texttt{DashboardLayout.tsx})~:
\begin{itemize}
    \item Bouton de rétraction de la sidebar
    \item Barre de recherche globale
    \item Bouton de notifications avec badge de compteur et panneau déroulant
    \item Menu utilisateur (DropdownMenu)~: nom du praticien, lien Profil, bouton Logout
\end{itemize}

% ============================================================
\section{Conclusion}

La phase de réalisation a permis de livrer une application fonctionnelle couvrant l'ensemble des besoins exprimés. Les fonctionnalités clés --- authentification sécurisée, gestion CRUD des patients et des rendez-vous, détection de conflits, calendrier hebdomadaire, créneaux disponibles et profil utilisateur --- sont pleinement opérationnelles.

L'architecture choisie (API REST + SPA React) offre une séparation claire des responsabilités et une expérience utilisateur fluide grâce au rendu côté client, au cache React Query et aux notifications toast instantanées.
